\section{Introduction}
\label{sec:intro}

Learning visuomotor policies from demonstrations involves a long chain
of software: dataset schemas, image encoding, normalisation statistics,
checkpointing, and an evaluation harness that must observe the world
exactly as training did. Each link can fail silently, and on a real
robot each failure is discovered only after hours or days of human
demonstration collection have already been spent. This paper describes
how we de-risk that campaign for our bimanual garment-manipulation rig:
we build the entire collect--train--evaluate loop inside the rig's
digital twin, drive it with scripted oracle demonstrators, and require
the loop to work end to end --- through three different policy families
--- before a single real demonstration is recorded.

The starting point is the platform described by our companion
teleoperation paper~\citep{so101teleop}: a pair of low-cost
five-degree-of-freedom arms on a printed rig, teleoperated from Meta
Quest hand controllers through a clutched retargeting stack. The rig has
a faithful digital twin: a single parametric CAD description generates
both the printed parts and the simulation model, so the simulated
geometry, camera poses and arm bases match the physical build by
construction.

A twin alone, however, is not a training environment. It needs an object
the grippers can actually hold under simulated contact; tasks whose
success is measurable; demonstrations whose distribution matches what
teleoperation will later produce; and an evaluation protocol that cannot
leak training scenarios. Building those four things --- and making them
reliable enough to train on --- is the subject of this paper.

Our contributions are:
\begin{itemize}
  \item a contact-manipulation extension of the digital twin: a graspable
    cube payload, fingertip contact pads, a bounded pinch force, and a
    timing design in which one control tick is an exact number of
    physics substeps (Section~\ref{sec:environment});
  \item two tasks with a seeded scenario generator and a disjoint
    train/validation/evaluation seed protocol, including a deliberate
    overfit-one-scenario mode that separates pipeline bugs from data
    problems (Section~\ref{sec:tasks});
  \item a scripted oracle that demonstrates through the \emph{full
    production teleoperation pipeline} by inverting the retargeting map,
    with closed-loop grasp and placement corrections standing in for the
    operator's eyes, plus a deterministic direct-drive fallback
    (Section~\ref{sec:oracles});
  \item an honest catalogue of every difficulty met while making
    contact-based demonstration collection reliable, and its resolution
    (Section~\ref{sec:difficulties});
  \item an experiment protocol --- oracle gate, gated collection,
    per-checkpoint validation, held-out evaluation --- whose results
    section grows with this living document
    (Section~\ref{sec:experiments}).
\end{itemize}

\textbf{TODO: add related work. Mention all the bimanual simulation benchmark in the literature. Like Bigym, Robosuite and etc.}

\textbf{TODO: also add the task of stacking cubes. Cubes can be everywhere but within the reach of the arms}

\textbf{TODO: do not add the platform rig in the simulation and this paper. Just fix the two arms side by side with 25cm apart. The third-person camera should have a front camera 45 degree from top to see the oepration space centre.}

\textbf{TODO: allow the arms to have wrist cameras but not with the 3D holder we designed. It should use the popular wrist camera type and holder online.}

\textbf{TODO: also add task of peg-in-hole. The peg is a graspable shape and the hole is a corresponding shape with a bigger size. The peg can be anywhere but within the reach of the arms. The hole is fixed in the centre of the operation space.}

\textbf{TODO: use white robot body, but keep the motor colours as black. Use dark gray table surface, whilst keeping the walls and rooftops with colours that are ascetically pleasing and it will be easy for vision-based policy to discern.}

\textbf{TODO: in the paper, clearly note down the extrinsics and instrinsics of the cameras}.

\textbf{TODO: add example pictures of each task and present them in one single big figure with labels.}

\textbf{TODO: have a subsection for each task to introduce the details of these tasks----including purpose, success criteria, reward function, evaluation metrics; also introduce their oracle and mock-teleopration oracle policies. Try to use pseudo-codes or state-flows to illustrate the oracle and mock-teleoperation oracle policies.}
